%% =============================================================================
%%  Como fazer um estudo temático (texto genérico: serve a qualquer tema)
%% =============================================================================
\chapter{Como fazer um estudo temático}

Um estudo temático parte de uma pergunta (\enquote{\temaPergunta}) e procura
a resposta em toda a Bíblia. O risco próprio desse tipo de estudo é juntar
versículos soltos que confirmam o que já se pensava. O roteiro abaixo existe
para evitar isso: reunir \emph{todos} os textos relevantes, ler cada um no
seu contexto, perceber como o assunto se desenvolve ao longo da Bíblia e só
então formular a síntese \parencite{osborne2006, duvall2012}.

\section{O roteiro em sete etapas}

\begin{ebtabela}{@{}>{\raggedright}p{.26\linewidth} >{\raggedright\arraybackslash}X@{}}
\EBcab{Etapa} & \EBcab{O que fazer}\\\midrule
1. A pergunta & formular a pergunta com precisão, dividi-la em subperguntas e registrar o que já penso sobre o assunto (os meus pressupostos)\\
2. Palavras e conceitos & levantar as palavras hebraicas e gregas do tema, com todas as ocorrências; procurar também sinônimos e ideias expressas sem essas palavras\\
3. Os textos & montar a lista de estudo, agrupar por assunto e marcar o peso de cada texto (central, apoio, difícil)\\
4. Leitura no contexto & estudar os textos centrais um a um: contexto, o que dizem, como contribuem para o tema\\
5. Desenvolvimento & ver como o tema aparece ao longo da Bíblia: o que muda, o que permanece, a quem cada orientação é dirigida\\
6. Diálogo & enfrentar os textos difíceis e ouvir como outras traduções, comentários e tradições os leem\\
7. Síntese & responder à pergunta com proposições numeradas, cada uma com a sua base textual; separar o que é claro do que fica em aberto; aplicar\\
\end{ebtabela}

\begin{alerta}[Palavra não é conceito]
Um tema pode estar presente sem a palavra que o nomeia, e a palavra pode
aparecer sem tratar do tema. Uma concordância é o ponto de partida, não a
lista final \parencite{barr1961, silva1994}. Em \enquote{Quem é Jesus?}, por
exemplo, o assunto está em títulos, ações, relações e declarações, não numa
única palavra; em \enquote{sangue}, a palavra está em quase todos os textos,
mas nem todos tratam do mesmo aspecto.
\end{alerta}

\begin{alerta}[Seleção de evidências]
Escolher só os textos que confirmam uma conclusão é o erro mais comum nos
estudos temáticos \parencite{carson1996}. Por isso a lista de ocorrências fica
no apêndice: ela mostra o que ainda não foi lido.
\end{alerta}

\section{Como usar este caderno}

\paragraph{Os dados.} A pasta \texttt{\temaDados} tem três arquivos de texto:
\texttt{textos.tsv} (a sua lista de estudo, que você escreve e reorganiza),
\texttt{vocabulario.tsv} (as palavras do tema) e \texttt{ocorrencias.tsv}
(gerado pela ferramenta). Os comandos estão no \texttt{MANUAL.md}:

\begin{nota}[Ferramentas]
\small\ttfamily
python3 ferramentas/tema.py ocorrencias H1818 G129 -o \temaDados\\
python3 ferramentas/tema.py vocabulario H1818 G129 -o \temaDados\\
python3 ferramentas/tema.py conferir \temaDados
\end{nota}

\paragraph{Grupos e peso.} Na coluna \texttt{grupo} de \texttt{textos.tsv}
você cria as suas categorias; o caderno agrupa os textos por elas, na ordem
em que aparecem no arquivo. A coluna \texttt{peso} aceita
\texttt{central}, \texttt{apoio} ou \texttt{dificil}; se ficar vazia, o caderno
impresso mostra caixinhas para decidir à mão
({\sffamily\tiny\color{suave}C\,\EBcaixa\;A\,\EBcaixa\;D\,\EBcaixa}).

\paragraph{Fichas de texto.} Para estudar um texto no contexto:

\begin{nota}[Exemplo]
\small\ttfamily
\textbackslash begin\{fichatexto\}[grupo]\{Lv 17:11\}\\
\ \ \textbackslash begin\{campo\}\{Contexto\} ... \textbackslash end\{campo\}\\
\ \ \textbackslash interlinear\{dados/biblia/lv017\_11.hbo.tsv\}\\
\textbackslash end\{fichatexto\}
\end{nota}

\paragraph{Síntese.} Cada conclusão vai numa \verb|\proposicao{afirmação}{textos}|.
Se um texto da base não aparece na lista de estudo, é sinal de que a lista
precisa ser revista.

\paragraph{Impresso ou digital.} No modo \texttt{impresso}
(\texttt{config.tex}), os campos que você ainda não preencheu aparecem com
linhas para escrever à mão; no modo \texttt{digital}, somem.
